%% =====================================================================
%%  Chapter 6: Conclusion
%%  Summarize the objective, method and results, and describe future work.
%%  Do not introduce new results here.
%% =====================================================================
\chapter{Conclusion}
\label{chap:conclusion}

\section{Summary}
\label{sec:summary}

\begin{writingnote}
Recall the objective and the problem stated in the Introduction, in the same words, and summarize what you did and what you found.
State the main results concretely, with the conditions and the key numbers, and make clear to what extent the objective was achieved.
Then explain what this study adds to the work reviewed in the Related Work chapter.
Do not generalize the conclusions beyond the conditions examined in the experiments.
For the typical structure of a conclusion and useful phrases, see~\cite{glasmandeal2010science,morleyphrasebank}.
\end{writingnote}

This thesis proposed \ldots{} to \ldots.
Computational experiments showed that the proposed method \ldots.

\section{Future Work}
\label{sec:future-work}

\begin{writingnote}
Derive future work from the limitations of this study that further research can address.
Take up limitations already stated in the Proposed Method or Computational Experiments chapter; do not raise new ones here for the first time.
Instead of a vague statement such as ``further improvement is needed'', describe what remains unsolved and how it could be addressed.
\end{writingnote}

%% Mark unfinished parts with \TODO{...}, e.g. \TODO{describe future work in detail}.
%% With the review option it is shown in red; without it, it is hidden and a warning is shown.
Future work includes \ldots.
